\section{Extraction: Two Failures With One Symptom}
\label{sec:extract}

\subsection{The containers}

\begin{table}[htbp]
\caption{Archive inventory, verified by the decompressor's own list command}
\label{tab:inventory}
\centering
\footnotesize
\begin{tabular}{@{}lrrrr@{}}
\toprule
\textbf{Box} & \textbf{Bytes} & \textbf{Files} & \textbf{Packed} & \textbf{Unpacked} \\
\midrule
\verb|fg-01| & 24{,}992{,}387{,}781 & 858   & 23{,}834\,MB & 51{,}419\,MB \\
\verb|fg-02| &  6{,}037{,}379{,}414 & 9{,}774 &  5{,}757\,MB &  5{,}833\,MB \\
\verb|fg-03| &  3{,}591{,}566{,}415 & 1     &  3{,}425\,MB &  3{,}750\,MB \\
\verb|fg-04| &  3{,}280{,}408{,}667 & 218   &  3{,}128\,MB &  4{,}011\,MB \\
\verb|fg-05| &    902{,}897{,}841 & 78    &    861\,MB &  1{,}393\,MB \\
\verb|fg-06| &        525{,}112 & 21    &      0\,MB &      2\,MB \\
\midrule
\textbf{Total} & \textbf{$\approx$37\,GB} & \textbf{10{,}950} & & \textbf{$\approx$64.8\,GB} \\
\bottomrule
\end{tabular}
\end{table}

Note the last row of Table~\ref{tab:inventory}. The smallest container, at
525\,KB and roughly one part in seventy thousand of the total, blocked the entire
project for two working sessions. It contains no game data whatsoever.

A reading trap in the decompressor's own output is worth recording: the size
fields reported by its list command are in \emph{megabytes} for the maximum, and
the running counter is a 32-bit signed integer that goes negative past 2\,GB.
Reading the counter rather than the maximum produces nonsense on any modern
archive.

\subsection{Failure 1: 0.0\% forever}

The installer, run under Wine, stopped at 0.0\% and stayed there. No error, no
timeout, no progress. The naive readings---``it is slow'', ``it deadlocked'',
``the GUI is broken''---are all wrong.

\subsubsection{Root cause}

The repack's external-compressor plugins are split across architectures: a 64-bit
helper process does the work, and a 32-bit plugin loaded inside the installer
communicates with it through shared memory. Tracing the relevant calls with
Wine's relay debugging revealed the mismatch:

\begin{itemize}
\item The 64-bit helper creates its mapping as
      \verb|Global\Read_MapFile_CLS-srep00|.
\item The 32-bit plugin opens \verb|Read_MapFile_CLS-srep00| --- \textbf{with no
      prefix}.
\end{itemize}

On Windows this works, and it works for a reason specific to the NT object
manager~\cite{ntobj}. A session's \verb|\Sessions\N\BaseNamedObjects| directory
is a \emph{shadow} directory: a name that misses there falls through to the
global \verb|\BaseNamedObjects|. \textbf{Wine implements no such fallback.} The
unprefixed open therefore misses, and the plugin responds the way a 2000s-era
Windows utility responds to an unexpected error---it raises a modal message box
reading \verb|Can't open read file mapping!| and waits.

That message box is never painted. The installer is not hung; it is waiting for a
click on a button that exists only in the window manager's data structures. This
is the first of two independent occurrences of this failure shape in this
project; the second, in the graphics stack, is Section~\ref{sec:deadlock}.

\subsubsection{Fix}

One byte per helper. Zeroing the leading \verb|G| of the \verb|"Global\"| string
literal truncates the prefix to the empty string, so both sides use the
session-local name and the open succeeds. Backups of the originals were retained.

\subsubsection{Why patching on disk is insufficient}

This is the part that generalises. Patching the helpers where they live has no
effect, because \textbf{the installer extracts its own fresh, unpatched copies of
every helper into a newly created temporary directory on each run}. A
hand-patched copy elsewhere is simply bypassed.

The working approach is a live patcher: a process that polls at 0.2-second
intervals and patches new copies as they appear. On launch it caught eight
unpatched helpers across two separate temporary directories within seconds.

The general lesson is that \emph{a self-extracting installer is a moving target,
and static patching of its payload is only valid if the payload is not
re-materialised at runtime}. The way to detect this class of problem is to notice
that the patch verifiably applied and verifiably had no effect---otherwise an
extremely confusing pair of facts.

\subsection{Failure 2: ``unsupported compression method''}

With the modal deadlock fixed, five of the six containers opened. The sixth---the
525\,KB one---failed differently:

\begin{lstlisting}
[file] work\work\build01.bat
ERROR: unsupported compression method
ERROR: archive data corrupted (decompression fails)
\end{lstlisting}

\subsubsection{Differential diagnosis: distinguishing the two failures}

The two failures present almost identically to a user: the process is alive and
nothing happens. They are distinguished by a single cheap observation, and
recording that observation is what prevented a second week on the first bug.

\begin{table}[htbp]
\caption{Telling the two extraction failures apart}
\label{tab:diff}
\centering
\footnotesize
\begin{tabularx}{\columnwidth}{@{}L{1.9cm}XX@{}}
\toprule
& \textbf{Failure 1 (modal)} & \textbf{Failure 2 (config)} \\
\midrule
CPU & 0\% & one core fully busy \\
Child process & helper alive & \textbf{none} \\
Bytes written & partial & \textbf{zero} \\
Time to fail & never & $\approx$2\,s \\
Root cause & namespace fallback & missing config file \\
\bottomrule
\end{tabularx}
\end{table}

\subsubsection{The wrong hypothesis, and how long it survived}

The error message is a direct statement that a compression method is not
supported, so the hypothesis was that the container used a codec the shipped
decompressor did not implement. Considerable effort went into it: an archive
scanner was written, the tail directory block was located at two byte offsets,
and a brute-force search over LZMA1 parameter space was implemented to recover
the per-file method table out of the compressed directory block.

None of it was the cause. The archive was correct the entire time.

\subsubsection{Root cause}

The decompression library reads its table of external compressors from a
configuration file, and it looks for that file at \textbf{exactly one hard-coded
absolute location}: the root of the system drive, \verb|c:\arc.ini|. Not its own
directory. Not the current working directory. Not the output directory. Not the
archive's directory.

This was established, rather than guessed, by tracing file operations with
\verb|WINEDEBUG=+file| from three different working directories. Every run
produced exactly one probe, always \verb|CreateFileW L"c:\arc.ini"|, and always
status \verb|c0000034| --- object name not found.

The file exists, in the toolchain directory. With the table absent, every
compression method that depends on an external tool resolves to ``unsupported,''
and extraction dies on the first file that uses one.

\subsubsection{Fix}

One line: copy the configuration file to the root of the emulated system drive.
The container that had blocked the project for two sessions opened on the first
attempt afterwards, as did every other container.

\subsection{Three signals that were available the whole time}

Each of the following was observable before any of the wasted effort began, and
each independently points away from the codec hypothesis.

\begin{enumerate}
\item \textbf{The list operation succeeded throughout.} The directory block
      parsed correctly and reported all files. \emph{A working list beside a
      failing extract points at the method table, not at the archive.} Were the
      archive corrupt or the format misunderstood, listing would fail first.
\item \textbf{Zero bytes written, failing in two seconds, is a parse-time
      failure.} A genuine codec or data fault writes a partial output file
      first, because it fails partway through a stream. Checking the destination
      directory---not the log---separates these instantly.
\item \textbf{``Unsupported compression method'' from this class of tool almost
      always means missing configuration, not a missing codec.} The method table
      is data, and data can simply be absent.
\end{enumerate}

\subsection{Dead ends, published so they are not repeated}

\begin{table}[htbp]
\caption{Failure 2: hypotheses tested and eliminated}
\label{tab:deadends}
\centering
\footnotesize
\begin{tabularx}{\columnwidth}{@{}L{2.7cm}X@{}}
\toprule
\textbf{Hypothesis} & \textbf{How it died} \\
\midrule
Corrupt configuration & Diffed against four copies; all identical \\
Destination filesystem & NTFS versus ext4 made no difference \\
Temporary path setting & Changing it made no difference \\
Regression from the Failure 1 patch & Helpers verified still patched \\
Unknown LZMA1 parameters & Brute-force search found nothing wrong, because
                            nothing was wrong \\
A Wine defect & Two different programs failed on the same file in the same way \\
\bottomrule
\end{tabularx}
\end{table}

\subsection{The standalone harness}

To separate ``the archive cannot be read'' from ``the installer cannot read the
archive,'' a minimal harness was built that loads the decompression library
directly and calls its extraction entry point, bypassing the installer framework
entirely.

The build is constrained in an instructive way. The library is 32-bit, and the
system's Wine is a new-WoW64 build with no 32-bit Unix support, so a 32-bit
\emph{Winelib} binary cannot be produced at all. What can be produced is a 32-bit
\emph{PE}: a freestanding executable with no C runtime and no Windows headers,
built with the LLVM toolchain already present.

\begin{lstlisting}
llvm-dlltool -m i386 -k -d kernel32.def -l kernel32.lib
clang --target=i386-pc-windows-msvc -nostdlib \
      -ffreestanding -O1 -c arcx32.c
lld-link /subsystem:console /entry:start \
      /nodefaultlib /machine:x86 arcx32.obj kernel32.lib
\end{lstlisting}

One implementation detail cost an afternoon and is a general hazard when calling
into legacy Windows libraries: \textbf{the progress callback must be declared}
\verb|__stdcall|. Declared \verb|cdecl|, the stack drifts by 16 bytes per
invocation and the program page-faults after a few hundred callbacks---far enough
into the run to look like a data-dependent fault rather than a calling-convention
error.
